\chapter{Corrections optiques}
\section{Activer les corrections disponibles}
Le groupe \emph{Corrections optiques} se trouve en fin de liste, avant Export.
Ses cases commandent la correction de l’objectif et la réduction de l’aberration
chromatique. Les capacités dépendent du RAW, des métadonnées et du profil
objectif disponible. Lire le diagnostic indiquant les corrections réellement
appliquées et l’éventuelle absence de profil LCP.

Le profil DCP du boîtier et le profil LCP de l’objectif ont des rôles différents.
Ne pas interpréter une case activée comme la preuve que toute correction
géométrique, de vignettage ou chromatique a effectivement pu être calculée.

\section{Comparer correctement}
Observer les lignes droites proches des bords, la luminosité des coins et
les franges sur les contours contrastés. Une correction de vignettage peut
éclaircir les bords et modifier l’impression lumineuse globale.
Pour comparer deux développements colorimétriques, garder les mêmes états
optiques ou les désactiver de manière cohérente des deux côtés.

\paragraph{Clavier.}
Utiliser \touche{Tab} pour atteindre une case puis \touche{Espace} pour la
basculer lorsqu’elle a le focus. Le même Espace n’active l’original temporaire
que dans la vue image. Il n’existe pas de raccourci global d’activation optique.
